\begin{problem}[One card from a deck]
A card is chosen at random from a standard deck. Find the probability that it is a king or a diamond.
\end{problem}

\begin{hint}
Use the addition rule and subtract the overlap only once.
\end{hint}

\begin{solution}
There are $4$ kings and $13$ diamonds, with one overlap: the king of diamonds.
\[
P(\text{king or diamond})=\frac{4+13-1}{52}=\frac{16}{52}=\frac4{13}.
\]
\end{solution}

\begin{problem}[Three coin tosses]
Three fair coins are tossed. Find the probability of getting exactly two heads.
\end{problem}

\begin{hint}
List the successful patterns or use the binomial idea with $n=3$ and $p=\frac12$.
\end{hint}

\begin{solution}
There are $\binom32=3$ ways to place two heads among three tosses.
\[
P(\text{exactly two heads})=\binom32\left(\frac12\right)^3=\frac38.
\]
\end{solution}

\begin{problem}[Without replacement]
A box contains $5$ white counters and $2$ black counters. Two counters are drawn without replacement. Find the probability that both are white.
\end{problem}

\begin{hint}
Multiply the changing fractions directly.
\end{hint}

\begin{solution}
\[
P(\text{both white})=\frac57\cdot\frac46=\frac{20}{42}=\frac{10}{21}.
\]
\end{solution}

\begin{problem}[Expected value from a table]
A random variable $X$ has
\[
P(X=1)=0.4,\qquad P(X=2)=0.1,\qquad P(X=4)=0.5.
\]
Find $E(X)$.
\end{problem}

\begin{hint}
Multiply each outcome by its probability and add.
\end{hint}

\begin{solution}
\[
E(X)=1(0.4)+2(0.1)+4(0.5)=0.4+0.2+2=2.6.
\]
\end{solution}

\begin{problem}[Hannah and the snack boxes]
Hannah may choose one item from a snack shop. The shop has:
\begin{itemize}[leftmargin=*]
  \item $6$ boxes of muffins, each muffin costing \$3.20
  \item $5$ boxes of cookies, each cookie costing \$2.40
  \item $4$ boxes of juice boxes, each juice box costing \$2.80
\end{itemize}

Hannah chooses a box at random, then takes one item from that box. Find the expected cost of the item she takes.
\end{problem}

\begin{hint}
First write down the probability of choosing each type of box. Then use a weighted average.
\end{hint}

\begin{solution}
There are
\[
6+5+4=15
\]
boxes in total.

So
\[
P(\text{muffin})=\frac6{15},\qquad
P(\text{cookie})=\frac5{15},\qquad
P(\text{juice box})=\frac4{15}.
\]

Let $X$ be the cost of the item chosen. Then
\[
E(X)=3.20\left(\frac6{15}\right)+2.40\left(\frac5{15}\right)+2.80\left(\frac4{15}\right).
\]

\[
E(X)=1.28+0.80+\frac{11.2}{15}=1.28+0.80+0.746\overline{6}=2.826\overline{6}.
\]

Therefore the expected cost is
\[
\$2.83 \text{ (approximately)}.
\]
\end{solution}

\begin{problem}[Reading probabilities from a table]
A random variable $X$ has the following probability distribution:
\[
\begin{array}{c|cccccc}
x & 0 & 1 & 2 & 3 & 4 & 5 \\
\hline
P(X=x) & 0.1 & 0.2 & 0.3 & 0.25 & 0.05 & 0.1
\end{array}
\]
Find:
\begin{enumerate}[(i)]
  \item $P(X \text{ is odd})$,
  \item $P(X\ge 2)$.
\end{enumerate}
\end{problem}

\begin{hint}
Add the probabilities of all outcomes that satisfy each condition.
\end{hint}

\begin{solution}
\textbf{(i)} The odd values are $1$, $3$, and $5$, so
\[
P(X \text{ is odd})=P(X=1)+P(X=3)+P(X=5)=0.2+0.25+0.1=0.55.
\]

\textbf{(ii)} The values at least $2$ are $2$, $3$, $4$, and $5$, so
\[
P(X\ge 2)=P(X=2)+P(X=3)+P(X=4)+P(X=5)=0.3+0.25+0.05+0.1=0.7.
\]
\end{solution}

\begin{problem}[Complement of an impossible pair]
Two fair dice are rolled. Find the probability that the sum is not $7$.
\end{problem}

\begin{hint}
There are six outcomes with sum $7$.
\end{hint}

\begin{solution}
\[
P(\text{sum }7)=\frac6{36}=\frac16.
\]
So
\[
P(\text{sum not }7)=1-\frac16=\frac56.
\]
\end{solution}

\begin{problem}[Single-stage conditional probability]
In a box there are $3$ green counters and $5$ yellow counters. One counter is drawn. Given that it is not yellow, find the probability that it is green.
\end{problem}

\begin{hint}
If it is not yellow, what possibilities remain?
\end{hint}

\begin{solution}
If the counter is not yellow, it must be green. Therefore
\[
P(\text{green}\mid \text{not yellow})=1.
\]
\end{solution}

\begin{problem}[Binomial single success]
A fair coin is tossed $5$ times. Find the probability of exactly one head.
\end{problem}

\begin{hint}
Use $\binom51\left(\frac12\right)^5$.
\end{hint}

\begin{solution}
\[
P(X=1)=\binom51\left(\frac12\right)^5=5\cdot \frac1{32}=\frac5{32}.
\]
\end{solution}

\begin{problem}[One coin and one die]
A fair coin and a fair die are used. Find the probability of getting a head and an even number.
\end{problem}

\begin{hint}
These two devices act independently, so multiply the probabilities.
\end{hint}

\begin{solution}
\[
P(\text{head and even})=\frac12\cdot\frac36=\frac12\cdot\frac12=\frac14.
\]
\end{solution}

\begin{problem}[No red counter]
A bag contains $3$ red counters and $4$ blue counters. One counter is drawn. Find the probability that it is not red.
\end{problem}

\begin{hint}
Either count blue directly or use the complement of red.
\end{hint}

\begin{solution}
\[
P(\text{not red})=1-\frac37=\frac47.
\]
\end{solution}

\begin{problem}[A letter from MIRANDA]
A letter is chosen at random from the word \texttt{MIRANDA}. Find the probability that it is:
\begin{enumerate}[(i)]
  \item an A,
  \item a vowel,
  \item a letter of the word \texttt{HACKING}.
\end{enumerate}
\end{problem}

\begin{hint}
There are $7$ letters in total. Count favourable letters carefully, including repeated letters when they appear in the word.
\end{hint}

\begin{solution}
The word \texttt{MIRANDA} has $7$ letters:
\[
\texttt{M, I, R, A, N, D, A}.
\]

\textbf{(i)} There are $2$ A's, so
\[
P(\text{A})=\frac27.
\]

\textbf{(ii)} The vowels are \texttt{I, A, A}, so there are $3$ vowels. Thus
\[
P(\text{vowel})=\frac37.
\]

\textbf{(iii)} The letters in \texttt{HACKING} are
\[
\texttt{H, A, C, K, I, N, G}.
\]
Among the letters of \texttt{MIRANDA}, the ones that belong to \texttt{HACKING} are
\[
\texttt{I, A, N, A},
\]
so there are $4$ favourable letters. Hence
\[
P(\text{a letter of \texttt{HACKING}})=\frac47.
\]
\end{solution}

\begin{problem}[HSC 2025-style: Bernoulli mean and variance]
A Bernoulli random variable $X$ has
\[
P(X=1)=\frac23,
\qquad
P(X=0)=\frac13.
\]
Find $E(X)$ and $\mathrm{Var}(X)$.
\end{problem}

\begin{hint}
For a Bernoulli random variable with success probability $p$, use $E(X)=p$ and $\mathrm{Var}(X)=p(1-p)$.
\end{hint}

\begin{solution}
Here $p=\frac23$, so
\[
E(X)=\frac23,
\qquad
\mathrm{Var}(X)=\frac23\cdot \frac13=\frac29.
\]
\end{solution}
